\documentclass[10pt,a4paper]{amsart}
\usepackage[margin=19mm]{geometry}
\usepackage{amsmath,amssymb,mathtools}
\usepackage[scr=rsfs,cal=euler]{mathalfa}
\usepackage{xfrac}
\usepackage[unicode,hidelinks,bookmarks=false]{hyperref}
\newcommand{\ie}{i.e.\ }
\newcommand{\cf}{cf.\ }
\newcommand{\resp}{respectively\ }
\newcommand{\Subsection}{\subsection}
\providecommand{\nobreakdash}{\nobreak\hbox{-}}
\setlength{\emergencystretch}{2em}
\multlinegap0pt
\usepackage[shortlabels]{enumitem}
\usepackage{amssymb}
\usepackage{caption}
\usepackage{algorithm}\usepackage{algpseudocode}
\usepackage{multirow}
\newtheorem{theorem}{Theorem}
\newtheorem{lemma}{Lemma}
\title{Compressed data structures for Heegaard splittings}
\author{Henrique Ennes, Cl\'ement Maria}
\date{Source: arXiv:2507.11406v2 (CC BY 4.0)}
\begin{document}
\maketitle

\noindent\emph{Source and licence.} Compressed data structures for Heegaard splittings. Version arXiv:2507.11406v2 (CC BY 4.0), distributed by arXiv under the Creative Commons Attribution 4.0 International licence, http://creativecommons.org/licenses/by/4.0/, as recorded in the arXiv licence metadata for this exact version and shown on the abstract page https://arxiv.org/abs/2507.11406v2. Full licence notice: \texttt{licenses/arxiv-2507.11406-CC-BY-4.0.txt}.\par\smallskip

\noindent\textit{Editorial scope.} Counted unit: the article's theorem maximal (source0.tex lines 991--993). The article's complete original proof of it is delivered with it (source0.tex lines 994--1023, one \texttt{proof} environment of 30 lines). No mathematics is written, repaired or re-derived: the statement and the proof are the article's own lines, in the article's own notation, and no retained line is substituted or re-flowed.  Retained uncounted as the source-local context the proof needs: lines 447--454; lines 633--635.  Nothing was omitted from any retained span, so the package carries no omission record.  No retained line contains a citation, so the wrapper carries no bibliography and no bibliography entry.  No retained line contains a figure environment, an image-inclusion command or drawing code, so no image is delivered and the package declares no asset dependency.  Editorial formatting only, in addition: an AMS \texttt{amsart} preamble that restates the article's own declarations and macros used by the retained lines, namely the environments \texttt{theorem}, \texttt{lemma} on separate counters, together with the standard packages those lines need.  The retained environments print in this document as Theorem 1, Lemma 1 in order of appearance, whereas the article prints the same items with the numbering of its own full document.  The complete native source of the article as distributed in the arXiv e-print for this exact version is bundled unchanged as \texttt{sources/181626/source0.tex}.
\begin{theorem}\label{th: erickson}
        Suppose that $(T,E_T(\alpha),N_T(\beta))$ is a genus $g$ Heegaard diagram of complexity $m$. In a unit-cost RAM model of bit-size $\log |\beta\cap E|=O(m/|T|)$, one may compute 
        \begin{enumerate}[label=(\alph*)]
            \item the street complex $S(T,\beta)$ in time $O(m|T|)$;
            \item SLPs of the crossing sequence of all streets of $S(T,\gamma)$ in time $O(m|T|)$ of \emph{total} complexity $O(m|T|)$;
             \item the SLPs $I_T(\beta)$ of complexity $O(m)$, in time $O(m)$.
        \end{enumerate}
\end{theorem}

\begin{theorem}[\texttt{DO A DISK SLIDE}]\label{th: do a disk slide}
    A diagram $(T,E_T(\alpha),N_T(\beta'))$, where $\beta'$ is the result of applying a disk slide on any two components of $\beta$, can be computed in time $O(m|T|)$.
\end{theorem}

\begin{theorem}[\texttt{EXTEND TO MAXIMAL}]\label{th: extend to maximal}
    Let $(T',E_{T'}(\alpha),N_{T;}(\beta))$ be a Heegaard diagram of complexity $m$. Then one can compute an equivalent maximal Heegaard diagram $(T,E_T(\alpha'),N_T(\beta'))$ in time $O(gm|T|)$.
\end{theorem}

\begin{proof}
    Without loss of generality, suppose $(T,E_T(\alpha),E_T(\beta))$ is a minimal diagram; we want to extend each multicurve to a pants decomposition. 
    Both cases can be treated similarly thanks to the following lemma.

\begin{lemma}\label{lm: get pants}
    Suppose $\gamma$ is an edged minimal system in a cellular embedding $T$ of a surface $\Sigma_g$. Then one can find a pants decomposition $\gamma'$, equivalent to $\gamma$, in time $O(g|T|)$.
\end{lemma}
\begin{proof}
    Note that $\Sigma_g\backslash \gamma$ is homeomorphic to the $2g$ punctured sphere, so let $d_1,\dots,d_{2g}$ represent the boundary components of this surface.
    We will find some $2g-3$ new separating curves $d_{2g+1},\dots, d_{3g-3}$ in $\Sigma_g$ that recursively define pairs of pants as follows.
    In step $i$, choose a curve $d_{g+j}, 0\leq j\leq i-1$ from a list of available curves. 
    Call this choice $s_1$ and pick a face $t$ adjacent to $s_1$. 
    Using breadth-first search, we find, in time $O(|T|)$, the shortest path $p$ in the dual graph starting from $t$ to any face adjacent to one of the available curves.
    Call the closest available curve $s_2$ and remove $s_1$ and $s_2$ from the list of available curves. 
    Note that we can encode $p$ as a list of faces of $T$.
    In time $O(\|E_T(s_1)\|)=O(|T|)$, we update $p$ by iteratively removing all faces adjacent to $s_1$ except for the last one.
    By construction, $p$ intersects each edge of the cellular complex at most once and does not intersect any face adjacent to an available curve, except for exactly one face adjacent to $E_T({s_1})$ and one face adjacent to $E_T({s_2})$.
    
    As in the proof of Theorem \ref{th: do a disk slide}, make $I_T(d_i)=I_T(s_2)^{-1}\circ I_T(p)^{-1}\circ I_T(s_1)\circ I_T(p)$, where $p$ is potentially trivial, and add $d_i$ to the list of available curves. 
    The curve $d_i$ separates a thrice-punctured sphere component with punctures defined by $s_1$, $s_2$, and $d_i$, from a $2g-i$ punctured sphere component, and which component is which can be identified in time $O(|T|)$.
    We then update $T$ to a cellular embedding of $\Sigma_0^{2g+i}$ by adding the edges and vertices defined by $I_T(d_i)$ to the graph and blocking paths in the dual graph transverse to them.
    While $T$ becomes more complicated at each step, because the arc $p$ is constructed to intersect a constant number of faces adjacent to the available curves and the curves $d_{g+j}$, $1\leq j\leq i-1$, are separating, at each step, breadth-first needs only to consider $O(|T|)$ many edges and vertices of the dual graph of $T$. 
    The algorithm terminates when $i=3g-3$, in which case we have defined, by construction, a maximal system $\gamma = \{c_1,\dots,c_g,d_{g+1},\dots,d_{3g-3}\}$.
\end{proof}

We now apply Lemma \ref{lm: get pants} to get some pants decomposition equivalent to $\alpha$, edged in a new triangulation $T'$. 
Because $|T'|=O(|T|)$, we can use the discussion of corner coordinates to find the coordinates $N_{T'}(\beta)$ in time $O(g|T|)$. 
We then compute some edged pants decomposition equivalent to $\beta$ in $S(T',\beta)$. 
Finally, by the very proof of Lemma \ref{lm: get pants} and part (b) of Theorem \ref{th: erickson}, this gives the normal coordinates of the new $\beta$ curves with respect to $T'$.
\end{proof}

\end{document}
